\documentclass[10pt,a4paper]{amsart}
\usepackage{amsmath,amssymb,amsthm,graphicx}
\usepackage{verbatim}
\newenvironment{reference}{\comment}{\endcomment}
%\newenvironment{reference}{}{}
\newenvironment{slogan}{\comment}{\endcomment}
\newenvironment{history}{\comment}{\endcomment}

% For commutative diagrams we use Xy-pic
\usepackage[all]{xy}

% We use 2cell for 2-commutative diagrams.
\xyoption{2cell}
% Theorem environments.
%
\theoremstyle{plain}
\newtheorem{theorem}[subsection]{Theorem}
\newtheorem{proposition}[subsection]{Proposition}
\newtheorem{lemma}[subsection]{Lemma}

\theoremstyle{definition}
\newtheorem{definition}[subsection]{Definition}
\newtheorem{example}[subsection]{Example}
\newtheorem{exercise}[subsection]{Exercise}
\newtheorem{situation}[subsection]{Situation}

\theoremstyle{remark}
\newtheorem{remark}[subsection]{Remark}
\newtheorem{remarks}[subsection]{Remarks}

\numberwithin{equation}{subsection}

% Macros
%
\def\lim{\mathop{\mathrm{lim}}\nolimits}
\def\colim{\mathop{\mathrm{colim}}\nolimits}
\def\Spec{\mathop{\mathrm{Spec}}}
\def\Hom{\mathop{\mathrm{Hom}}\nolimits}
\def\Ext{\mathop{\mathrm{Ext}}\nolimits}
\def\SheafHom{\mathop{\mathcal{H}\!\mathit{om}}\nolimits}
\def\SheafExt{\mathop{\mathcal{E}\!\mathit{xt}}\nolimits}
\def\Sch{\mathit{Sch}}
\def\Mor{\mathop{\mathrm{Mor}}\nolimits}
\def\Ob{\mathop{\mathrm{Ob}}\nolimits}
\def\Sh{\mathop{\mathit{Sh}}\nolimits}
\def\NL{\mathop{N\!L}\nolimits}
\def\CH{\mathop{\mathrm{CH}}\nolimits}
\def\proetale{{pro\text{-}\acute{e}tale}}
\def\etale{{\acute{e}tale}}
\def\QCoh{\mathit{QCoh}}
\def\Ker{\mathop{\mathrm{Ker}}}
\def\Im{\mathop{\mathrm{Im}}}
\def\Coker{\mathop{\mathrm{Coker}}}
\def\Coim{\mathop{\mathrm{Coim}}}

% Boxtimes
%
\DeclareMathSymbol{\boxtimes}{\mathbin}{AMSa}{"02}

%
% Macros for moduli stacks/spaces
%
\def\QCohstack{\mathcal{QC}\!\mathit{oh}}
\def\Cohstack{\mathcal{C}\!\mathit{oh}}
\def\Spacesstack{\mathcal{S}\!\mathit{paces}}
\def\Quotfunctor{\mathrm{Quot}}
\def\Hilbfunctor{\mathrm{Hilb}}
\def\Curvesstack{\mathcal{C}\!\mathit{urves}}
\def\Polarizedstack{\mathcal{P}\!\mathit{olarized}}
\def\Complexesstack{\mathcal{C}\!\mathit{omplexes}}
% \Pic is the operator that assigns to X its picard group, usage \Pic(X)
% \Picardstack_{X/B} denotes the Picard stack of X over B
% \Picardfunctor_{X/B} denotes the Picard functor of X over B
\def\Pic{\mathop{\mathrm{Pic}}\nolimits}
\def\Picardstack{\mathcal{P}\!\mathit{ic}}
\def\Picardfunctor{\mathrm{Pic}}
\def\Deformationcategory{\mathcal{D}\!\mathit{ef}}
\usepackage{cleveref}
\crefname{originallemma}{original lemma}{original lemmas}
\Crefname{originallemma}{Original lemma}{Original lemmas}
\makeatletter
\newlabel{lemma-duality-proper-over-field}{{0FVV}{source}{}{sourcepointer-1}{}}
\newlabel{lemma-duality-proper-over-field@cref}{{[originallemma][0][]0FVV}{[1][0][]source}{}{}{}}
\makeatother
\crefname{originalsection}{original section}{original sections}
\Crefname{originalsection}{Original section}{Original sections}
\makeatletter
\newlabel{section-dualizing-module}{{0AWH}{source}{}{sourcepointer-2}{}}
\newlabel{section-dualizing-module@cref}{{[originalsection][0][]0AWH}{[1][0][]source}{}{}{}}
\makeatother
\crefname{originallemma}{original lemma}{original lemmas}
\Crefname{originallemma}{Original lemma}{Original lemmas}
\makeatletter
\newlabel{lemma-dualizing-module-CM-scheme}{{0AWT}{source}{}{sourcepointer-3}{}}
\newlabel{lemma-dualizing-module-CM-scheme@cref}{{[originallemma][0][]0AWT}{[1][0][]source}{}{}{}}
\makeatother
\title{The dualizing module of a Cohen-Macaulay proper scheme over a field and its Serre duality pairings}
\author{The Stacks Project Authors}
\date{stacks-project:duality.tex:local-snapshot-20261012; GNU FDL 1.2 or later}
\begin{document}
\maketitle
\noindent\textit{Editorial note.} All delivered mathematical text is the author's own, copied from the pinned Stacks Project chapter file duality.tex; no statement or proof step was rewritten, reordered or removed.\par\medskip
\noindent\textit{Editorial note.} Preamble provenance: the pinned chapter file opens with an input of the shared preamble and carries no preamble text of its own, so the project's single Stacks-wide preamble staged at sources/stacks-duality/preamble.tex is excerpted verbatim under the recipe fixed by LANE\_BRIEF 22.1, as three spans, lines 8-13 (verbatim package and the reference/slogan/history environments), lines 14-18 (Xy-pic setup) and lines 154-217 (the native theorem declarations and mathematical macros). Nothing was added, removed or redefined; the wrapper class and the amsmath/amssymb/amsthm/graphicx base packages are supplied by the assembler.\par\medskip
\noindent\textit{Editorial note.} The counted statement and proof contain no cross reference, no citation and no included file, so the delivered body is self contained.\par\medskip
\noindent\textit{Editorial note.} Canonical location of the counted proposition: Stacks tag 0FVZ (sources/stacks-duality/pinned-original-source.tex line 7106).\par\medskip
\noindent\textit{Allowed standard prerequisites (editorial).}\begin{itemize}
\item Duality for a proper scheme over a field. Same chapter, Lemma labelled lemma-duality-proper-over-field, cited at proof lines 7123-7131. Supplies the dualizing complex, the pairing with H\textasciicircum{}0 and the duality isomorphisms.
\item The dualizing module of a Cohen-Macaulay scheme is Cohen-Macaulay. Same chapter, Lemma labelled lemma-dualizing-module-CM-scheme, cited at proof lines 7127-7129. Gives the Cohen-Macaulay property of omega\_X and its support.
\end{itemize}\medskip
\ifdefined\hypertarget\hypertarget{sourcepointer-1}{}\fi
\noindent\textit{Source reference.} Original lemma 0FVV: sources/stacks-duality/pinned-original-source.tex, source lines 6934--6934. Same-chapter pointer into the author environment lemma at source line 6934; invoked at proof line 7124 of the counted proof. It is retained only as a navigation pointer; every load-bearing step is inside the delivered span.\par\medskip
\ifdefined\hypertarget\hypertarget{sourcepointer-2}{}\fi
\noindent\textit{Source reference.} Original section 0AWH: sources/stacks-duality/pinned-original-source.tex, source lines 6031--6031. Same-chapter pointer into the author environment lemma at source line 6031; invoked at proof line None of the counted proof. It is retained only as a navigation pointer; every load-bearing step is inside the delivered span.\par\medskip
\ifdefined\hypertarget\hypertarget{sourcepointer-3}{}\fi
\noindent\textit{Source reference.} Original lemma 0AWT: sources/stacks-duality/pinned-original-source.tex, source lines 6214--6214. Same-chapter pointer into the author environment lemma at source line 6214; invoked at proof line 7128 of the counted proof. It is retained only as a navigation pointer; every load-bearing step is inside the delivered span.\par\medskip
\begin{lemma}
\label{lemma-duality-proper-over-field-CM}
Let $X$ be a proper scheme over a field $k$ which is Cohen-Macaulay
and equidimensional of dimension $d$. The module $\omega_X$
of Lemma \ref{lemma-duality-proper-over-field} has the following properties
\begin{enumerate}
\item $\omega_X$ is a dualizing module on $X$
(Section \ref{section-dualizing-module}),
\item $\omega_X$ is a coherent Cohen-Macaulay module whose support is $X$,
\item there are functorial isomorphisms
$\Ext^i_X(K, \omega_X[d]) = \Hom_k(H^{-i}(X, K), k)$
compatible with shifts and distinguished triangles for $K \in D_\QCoh(X)$,
\item there are functorial isomorphisms
$\Ext^{d - i}(\mathcal{F}, \omega_X) = \Hom_k(H^i(X, \mathcal{F}), k)$
for $\mathcal{F}$ quasi-coherent on $X$.
\end{enumerate}
\end{lemma}

\begin{proof}
It is clear from Lemma \ref{lemma-duality-proper-over-field}
that $\omega_X$ is a dualizing module (as it is the left most
nonvanishing cohomology sheaf of a dualizing complex).
We have $\omega_X^\bullet = \omega_X[d]$ and $\omega_X$ is Cohen-Macaulay
as $X$ is Cohen-Macualay, see Lemma \ref{lemma-dualizing-module-CM-scheme}.
The other statements follow from this combined with the
corresponding statements of Lemma \ref{lemma-duality-proper-over-field}.
\end{proof}

\end{document}
